% Imporditud: tools/import_eksperimendid.py
% Allikas: Eesti füüsikaolümpiaadi eksperimendi ülesannete
%   kogu (Rael Kalda, 2023), ülesanne Ü169, lahendus L169.
\setAuthor{EFO žürii}
\setRound{piirkonnavoor}
\setYear{2010}
\setNumber{G E2}
\setDifficulty{5}
\setTopic{geomeetriline optika}
\setVahendid{briljandikujuline läbipaistev ese, millimeeterpaber, joonlaud, pliiats. Eseme tipu moodustavate vastastahkude vaheline nurk on 90$^\circ$. Eseme kõrgus, st tipu kaugus suurimast tahust on h = 32 mm. Õhu murdumisnäitaja võib lugeda võrdseks ühega}
\setPunktid{14}
\setAllikas{Eksperimendi kogu Ü169, lahendus lk 127}
\setLahendusseis{toores}

\prob{Briljant}
Kasutades täieliku sisepeegeldumise efekti, leidke optilise detaili murdumisnäitaja.
\textit{Märkus:} võib osutuda kasulikuks teadmine, et kehtib valguse murdumise seadus $n_1$ sin $\alpha_1$ = $n_2$ sin $\alpha_2$, kus $n_1$ ja $n_2$ on keskkondade murdumisnäitajad ning $\alpha_1$ ja $\alpha_2$ on nurk pinnanormaali ja valguskiire vahel vastavas keskkonnas.

\hint

\solu
Lahendus 1.

Märgime paberil kaks punkti vahekaugusega nt 10 mm. Ühe juurde (P ) asetame

briljandi tipu, teise (punkti S) katame aga ühe külgtahuga. Asetame joonlaua bril-

jandi suurimale tahule, tihedalt vastu tahku nii, et joonlaua üks ots toetub paberile --- punktis Q, briljandi tipust P kaugusel 32 mm/ 2 = 45 mm, vt joonist. Nüüd

liigutame pead üles-alla, otsides üles koha, kus läbi suurima tahu vaadates kaob

külgtahuga kaetud punkt sisepeegelduse tõttu. Võtame sellel kadumise hetkel lu-

gemi punkti kujutisega kohakuti olevast joonlaua punktist R, saades teada vahe-

maa QR = 26 mm. Sisepeegelduv kiir on seega SR. Et QS = 35 mm, siis koosi-

nusteoreemist SR = 24,8 mm, st siinusteoreemist sin $\angle$QSR = sin 45$^\circ$ $\cdot$ 26 ning 24,8

n = 1/ cos $\angle$QSR $\approx$ 1,49.

Lahendus 2. Kui vaadata läbi eseme põhjatahu mingit valget pinda nii, et alguses on ese silmale lähedal ja hakata silmast eemale nihutama, on võimalik leida täielikud sisepeegeldused tumedate aladena üheaegselt kõigilt kaheksalt tahult. Sisepeegeldustest moodustatud hulknurk on seda suurem mida kaugemale eseme nihutame ja kaugusel umbes 210 mm silmast täidab ta kogu põhjatahu ala, mille maksimaalne läbimõõt on 33 mm. Kaugemale liikudes on servmised kiired erinevalt lähedasest olekust järjest rohkem risti otsatahu pinnaga ja mingist nurgast alates leiab aset täielik sisepeegeldumine. Et ese on sümmeetriline, teeme joonise kiirte käigu kohta piirjuhul.

Jooniselt näeme, et täieliku sisepeegelduse tekkimise tingimuseks on 1/n = sin(45$^\circ- \alpha'$) = (cos $\alpha'-$sin $\alpha'$)/ 2. Nüüd võime arvestada, et n$\alpha$u$'$ .rkN$\alpha$ii$'$soiins,v1ä/inke=ja(s1e$-$eg$\alpha$a $'$v)õ/im2e. kasutada ligikaudseid avaldisi cos $\alpha' \approx$ 1 ja sin $\alpha'$ $\approx$

Mniunrgdnum=iss2ea+du$\alpha$se$\approx$sts2aa+m12e61,a05s$\approx$en1d,a4d9a. sin $\alpha'$ = sin $\alpha$/n, st $\alpha' \approx \alpha$/n : 2 = n $- \alpha$

Märkus: Eelpooltoodud ligikaudse arvutuse asemel on võimalik läbi viia ka täpseid

trigonomeetrilisi arvutusi, kuid vastus muutub sellest vähem kui 0,003 võrra.

Lahendus 3. (Ligikaudne, 9 punkti). Märgime paberile punkti. Asetame briljandi külgtahuga punkti P peale. Vaatame läbi briljandi aluse (suurima tahu) punktile P . Paneme tähele, et punkt kaob ära siis, kui vaadata alusele üsna täpselt ristisuunas. Punkti P juures oleval tahul toimub täielik sisepeegeldus, läbi aluse läheb valgus murdumiseta. Ülesande tingimustest on nurk aluse ja külgtahu vahel 45 kraadi. Sisepeegelduse piirtingimusest sin $\alpha$ = 1/n, kus $\alpha$ = 45 kraadi, leiame n = 1/ sin $\alpha$ = 2 $\approx$ 1,41.

Märkus 1: See lahendus läheb läbi ka juhul kui asetada briljant algselt suuremale tahule ja vaadata läbi külgtahu, kuid sel juhul ei ole külgtahu ja vaatesuuna ligikaudse ristiolemise asjaolu nii läbinähtav.

Märkus 2: Seda meetodit on võimalik edasi arendada, tehes lahendusega 2 analoogseid mõõtmisi ja arvutusi. Kui see analüüs on läbi viidud korrektselt, siis on lahendus väärt täispunkte.
\probend
